\documentclass[12pt, oneside]{book}
% -------------------
% PACKAGES
% Basic font setup
\usepackage[english]{babel}
\usepackage[utf8]{inputenc}
\usepackage[T1]{fontenc}
\usepackage{lmodern}

% Figures
\usepackage{graphicx}
\usepackage[font={footnotesize, sf}, labelfont=bf]{caption} 
\graphicspath{{figs/}{../figs/}}

% Maths tools
\usepackage{amsmath, amssymb, mathtools}
\usepackage{amsthm, thmtools}
\usepackage{tikz-cd}
\usepackage{siunitx}

\sisetup{detect-family = true}

% TOC Format
% \usepackage[titles]{tocloft}
% \renewcommand{\cftchappagefont}{\sffamily}   
% \renewcommand{\cftsecpagefont}{\sffamily}   
% \renewcommand{\cftsubsecpagefont}{\sffamily}   

% References
\usepackage{csquotes}

% Utilites
\usepackage{enumerate}

% -------------------
% CUSTOMISATION
% Geometry setup
\usepackage{geometry}
\usepackage{setspace}
\geometry{
	a4paper,
	right = 2.5cm,
	left = 2.5cm,
	bottom = 3cm,
	top = 3cm
}
\renewcommand{\baselinestretch}{1.3}

% -------------------
% Reference setup
\usepackage{hyperref}
\usepackage[english, capitalise, noabbrev]{cleveref}
\hypersetup{
	colorlinks,
	linkcolor = {red!50!blue},
	citecolor = {red!50!blue},
	urlcolor = {red!50!blue},
	linktoc = page
}

% -------------------
% BIBLIOGRAPHY
\usepackage[style=alphabetic, citestyle=alphabetic]{biblatex}
\addbibresource{refs.bib}

% -------------------
% Theorem environments
\newcommand{\qedtriangle}{\ensuremath{\triangle}}
\newcommand{\qedtriangledown}{\ensuremath{\bigtriangledown}}
\allowdisplaybreaks

\declaretheoremstyle[spaceabove=6pt, spacebelow=6pt, headfont=\bfseries,
notefont=\normalfont, notebraces={(}{)}, qed=\qedtriangle]{definition}
\declaretheoremstyle[spaceabove=6pt, spacebelow=6pt, headfont=\bfseries,
notefont=\normalfont, notebraces={(}{)}, qed=\qedtriangledown]{example}

\declaretheorem[name=Theorem, refname={theorem,theorems}, Refname={Theorem,Theorem},
numberwithin = chapter ]{theorem}

\declaretheorem[name=Proposition, refname={proposition,propositions},
Refname={Proposition,Propositions}, numberlike=theorem]{proposition}

\declaretheorem[name=Lemma, refname={lemma,lemmas},
Refname={Lemma,Lemmas}, numberlike=theorem]{lemma}

\declaretheorem[name=Definition, style=definition, refname={definition,definitions},
Refname={Definitio,Definitions}, numberlike=theorem]{definition}

\declaretheorem[name=Example, style=example, refname={example,examples},
Refname={Example,Examples}, numberlike=theorem]{example}

\declaretheorem[name=Remark, style=remark, refname={remark,remarks},
Refname={Remark,Remarks}, numberlike=theorem]{remark}

% -------------------
% LAYOUT
\usepackage[bf,sf,small,pagestyles]{titlesec}
\usepackage{titling}

% Pagestyle defintions
\newpagestyle{main}[\sffamily \footnotesize]{
	\sethead*{\ifthesection{{\bfseries \thesection} \sectiontitle}{}}{}{{\bfseries Chapter \thechapter.} \chaptertitle}
	\headrule
	\footrule
	\setfoot*{}{}{\thepage}
}
\renewpagestyle{plain}[\sffamily \footnotesize]{
	\footrule
	\setfoot*{}{}{\thepage}
}
\renewpagestyle{empty}{}

% Format of chapter titles
\titleformat{\chapter}[block]{\sffamily \bfseries \Huge}{\filleft \large Chapter \Huge \thechapter\\}{0pt}{\Huge \titlerule[1pt] \vspace{1ex} \filleft}

% -------------------
% SUBFILES
\usepackage{subfiles}

\newcommand{\set}[1]{\left\{ #1 \right\}}
\newcommand{\abs}[1]{\left\lvert #1 \right\rvert}

\newcommand{\N}{\mathbb{N}}
\newcommand{\Z}{\mathbb{Z}}
\newcommand{\R}{\mathbb{R}}
\newcommand{\F}{\mathbb{F}}
\renewcommand{\L}{\mathcal{L}}

\DeclareMathOperator{\Sym}{Sym}
\DeclareMathOperator{\Lor}{Lor}
\newcommand{\ext}{\mathsf{\Lambda}}
\newcommand{\ud}{\,\mathrm{d}}
\newcommand{\id}{\,\mathrm{id}}
\renewcommand{\d}{\mathrm{d}}

\newcommand{\iprod}{\mathbin{\lrcorner}}

\newcommand{\Top}{\mathsf{Top}}
\newcommand{\Grp}{\mathsf{Grp}}
\newcommand{\Ab}{\mathsf{Ab}}
\newcommand{\Vect}{\mathsf{Vect}}
\newcommand{\Ring}{\mathsf{Ring}}
\newcommand{\Simp}{\mathsf{Simp}}
\newcommand{\Ch}{\mathsf{Ch}}


\makeatletter
\newcommand*{\defeq}{\mathrel{\rlap{%
    \raisebox{0.3ex}{$\m@th\cdot$}}%
  \raisebox{-0.3ex}{$\m@th\cdot$}}%
	=
}
\makeatother

\DeclareMathOperator{\im}{im}

% Category theory
\newcommand{\longto}{\longrightarrow}
\newcommand{\into}{\hookrightarrow}

\newcommand{\MKNN}{M\( k \)NN }
\newcommand{\hombound}[2]{\bar{\partial}_{#1}^{#2}}
\newcommand{\hbound}[1]{\bar{\partial}^{#1}}

% Cycles and boundaries
\newcommand{\cyc}[1]{\ker{\partial}_{#1}}
\newcommand{\bound}[1]{\im{\partial}_{#1}}

\newcommand{\pcyc}[2]{\ker{\partial}_{#1}^{#2}}
\newcommand{\pbound}[2]{\im{\partial}_{#1}^{#2}}


\title{A variational derivation of the field equations of an action-dependent
	Einstein-Hilbert Lagrangian}
\author{Arnau Mas}
\date{June 2021}

\begin{document}
\begin{titlepage}
	\setstretch{2.0}
	\centering \sffamily

	\vspace*{2cm}

	\includegraphics[width = 8cm]{logo-uab}

	\vspace{2cm}

	{\Large \itshape Physics BSc. Undergraduate Thesis} \\
	{\large \itshape June 2021}

	\vspace{10pt}
	\hrule
	\vspace{10pt}
	{\bfseries \Large A variational derivation of the field equations of an action-dependent
	Einstein-Hilbert Lagrangian}

	\vspace{10pt}
	\hrule		
	\vspace{2cm}

	{\Large Arnau Mas Dorca}

	\vspace{1cm}
	{\large \itshape Supervised by}

	{\Large Dr Jordi Gaset}
\end{titlepage}

\selectlanguage{english}
\pagestyle{plain}
\frontmatter
\hspace{0pt}
\vfill

\begin{itshape}
	I want to express deep gratitude to my supervisor Dr Jordi Gaset. This thesis would not
	have been possible without his generosity with his time and willingness to accomodate
	my interests. Having weekly discussion meetings has been a very rewarding experience and
	has made the process of research that much more engaging. 

	I also wish to mention Adrià Marín, who has also been working under Dr Gaset on a
	related topic and we have had many a productive discussion on our theses. 

	Of course I must also manifest my gratitude to my family for their support and
	understanding. This work is the capstone to five years of study, the last of which
	having taken place in the midst of a global pandemic, and I am sure they will be very
	proud of it. 
\end{itshape}

\vfill
\hspace{0pt}

\pagebreak
\section*{Abstract}
We derive the equations of motion of an action-dependent version of the Einstein-Hilbert Lagrangian, as a specific instance of the Herglotz variational problem. Action-dependent Lagrangians lead to dissipative dynamics, which cannot be obtained with the standard method of Lagrangian field theory. First-order theories of this kind are relatively well understood, but examples of singular or higher-order action-dependent field theories are scarce. This work constitutes an example of such a theory. By casting the problem in clear geometric terms we are able to obtain a Lorentz invariant set of equations, which contrasts with previous attempts. 


	{\small \sffamily \tableofcontents}

\mainmatter
\pagestyle{main}

\chapter{Background and motivation}
\subfile{chapters/chapter-1}

\chapter{The Herglotz variational problem}\label{ch:herglotz}
\subfile{chapters/chapter-2}

\chapter{Action-dependent Einstein gravity}\label{ch:einstein}
\subfile{chapters/chapter-3}

\chapter{Significance of the equations}\label{ch:significance}
\subfile{chapters/chapter-4}

\chapter{Conclusions}\label{ch:conclusions}
\subfile{chapters/chapter-5}

\backmatter

\pagestyle{plain}
\setquotestyle{english}
\printbibliography

\end{document}
